\section{Einleitung}
\label{sec:einleitung}
 
% --- 1.1 Motivation ---
\subsection{Motivation und Problemstellung}
\label{subsec:motivation}
 
Die Klassifikation von Fernerkundungsbildern beschäftigt sich mit der Frage, welche Landnutzungs- oder Landbedeckungsklasse eine gegebene Bildszene repräsentiert.
Dies ist von hoher praktischer Relevanz, da Satelliten- und Luftbilder in immer größerem Umfang verfügbar werden und eine manuelle Auswertung dieser Datenmengen weder zeitlich noch wirtschaftlich realisierbar ist.
Automatisierte Verfahren sind daher unerlässlich, um beispielsweise Veränderungen in der Landnutzung zu erkennen, landwirtschaftliche Flächen zu überwachen oder Katastrophenschäden zu identifizieren.
 
% --- 1.2 Zielsetzung ---
\subsection{Zielsetzung der Arbeit}
\label{subsec:zielsetzung}
 
Ziel dieser Arbeit ist es, eine eigene CNN-Architektur zu entwerfen und zu trainieren, die Bilder des PatternNet-Datensatzes möglichst genau in die 38 vorgegebenen Szenenklassen klassifiziert.
Im Vordergrund steht dabei nicht die Anwendung einer vortrainierten Standardarchitektur, sondern der eigenständige Entwurfsprozess: die systematische Auswahl von Schichten und anderen Parametern sowie die Analyse ihres Einflusses auf die Klassifikationsleistung.

% --- 1.3 Aufbau der Arbeit ---
\subsection{Aufbau der Arbeit}
\label{subsec:aufbau}
 
Der weitere Bericht ist wie folgt gegliedert:
Abschnitt~\ref{sec:methode} beschreibt den verwendeten Datensatz, die Vorverarbeitung, die entwickelte Netzarchitektur sowie das Trainingsverfahren.
In Abschnitt~\ref{sec:ergebnisse} werden die erzielten Ergebnisse anhand der gewählten Evaluationsmetriken präsentiert.
Abschnitt~\ref{sec:diskussion} diskutiert die Ergebnisse und identifiziert Stärken sowie Limitationen des gewählten Ansatzes.
Abschnitt~\ref{sec:fazit} fasst die Arbeit zusammen und gibt einen Ausblick auf mögliche Weiterentwicklungen.